% !TEX TS-program = xelatex
% !TEX encoding = UTF-8 Unicode
% !Mode:: "TeX:UTF-8"

\documentclass{resume}
\usepackage{zh_CN-Adobefonts_external} % Simplified Chinese Support using external fonts (./fonts/zh_CN-Adobe/)
%\usepackage{zh_CN-Adobefonts_internal} % Simplified Chinese Support using system fonts
\usepackage{linespacing_fix} % disable extra space before next section
\usepackage{cite}

\begin{document}
\pagenumbering{gobble} % suppress displaying page number

\name{张志锐(ZHIRUI ZHANG)}
% \name{越空}

\basicInfo{
  \email{zrustc11@gmail.com} \textperiodcentered\ 
  \phone{(+86) 188-106-36928} \textperiodcentered\ 
  \text{\faFilesO\ https://zrustc.github.io/}
  }

\section{\faGraduationCap\  个人简介}

\textit{在深度学习领域具备丰富研究经验和落地实践，对世界模型以及AI原生应用/硬件方向兴趣浓厚；对AGI/ASI有强烈信仰，并推演出第一性原理；熟悉文本/多模态基础大模型预训练范式和后训练范式，曾深度参与近万亿级别MoE基础大模型项目（全流程），同时具备丰富业务大模型落地经验（包括终端和云侧）；} \\

\textit{中国科学技术大学本博，MSRA-USTC联合培养博士，2019年博士毕业，导师是沈向洋和陈恩红，副导师是李沐，博士期间主要研究方向为深度学习和NLP，Google Scholar Citations达到2900+，h-index为24，发表过40篇以上顶会论文。现任华为终端20级技术专家，负责小艺大模型通用后训练以及推理模型，支撑多个场景大模型能力落地。先后曾在微软研究院（中国和美国）实习和阿里达摩院（阿里星）、腾讯AI Lab和阶跃星辰（初创）工作过，主导阿里通用翻译和腾讯翻译构建，提出业界首个个性化机器翻译、IWSLT2023端到端语音到语音翻译第一名、2018年机器翻译技术报告主要贡献者（新闻领域实现人类可比水平）；}

% \datedsubsection{\textbf{中国科学技术大学}}{2014 -- 2019}
% \textit{中国科学技术大学-微软亚洲研究院联合培养博士}\ \ 计算机科学与技术 \\
% \textit{导师: 沈向洋、陈恩红 \ \ 副导师：李沐} \\
% \textit{研究方向：深度学习、自然语言处理、机器翻译、自然语言生成} \\
% \text{Google Scholar Citations: $2400^{+}$,\ h-index: 20+}
% \datedsubsection{\textbf{中国科学技术大学}}{2010 -- 2014}
% \textit{本科}\ \ 计算机科学与技术


\textit{}

\section{\faUsers\ 主要经历}

\datedsubsection{\textbf{华为}，深圳/东莞}{2024.3 -- 2025.9}
\role{终端-小艺业务部，技术专家A (20级)，从零组建团队（团队规模8人/1名天少），1次A绩效}
\text{负责小艺大模型通用后训练、推理模型以及大模型场景落地}
\begin{itemize}
\item 通用后训练 (2024.9-2025.9)：(1) 负责小艺大模型通用后训练，梳理盘古大模型通用后训练流程和问题，牵引诺亚实现SFT3.0构建，支撑通用数据和模型高效迭代；(2) Qwen2.5-7B基座模型+该训练策略，\textbf{在模型Release一个月内和资源限制条件下}实现主流开源测试集（知识、推理、数学、代码、指令遵从、工具调用等）均超过Qwen2.5-7B-Instruct，其中AlpacaEval 2.0和Arena-hard提升10\%+, 内部测试集各个场景提升5\%--10\%；相同策略在盘古模型38B/135B-V3上实现10\%+相对分位值提升；(3) 相应策略应用于多模态大模型，基于Qwen3-8B构建多模态大模型+后训练高质量指令数据退火+SFT3.0训练策略，实现Qwen3-8B-VL的通用能力持平/超过Qwen2.5-32B-VL（数学和学科类持平、指令遵从+5\%、ArenaHard持平等）；(4) \textbf{支撑小艺对话Live态430亮点特性发布和上线（基于DeepSeek-V3场景微调）}；RL+GenRM训练策略显著提升小艺Live态对话满意度（+23\%，满意度3.35->4.15）；
\item 慢思考推理模型 (2024.9-2025.9)：(1) 负责通用慢思考推理模型构建以及前沿技术探索，完成多项技术探索和验证，包括BoN+PRM能力边界、快慢思考融合和可控思考长度策略（早于Qwen3发布）、模型解题与判别能力的双生特性、基于置信度的RL优化算法、环境交互端到端RL训练 (Agentic RL)等，\textbf{推演出推理模型的第一性原理(通向AGI/ASI)并完成OpenAI的O系列low/medium/high/pro思考模式的技术方案解密}；(2) 小艺深度解题智能体技术负责人，落地面向AI教育场景的快慢思考融合推理模型、AI互动讲解亮点特性打造等，\textbf{在2025电信打榜实现AI学习能力第一}；(3) 小艺数字人预研项目负责人，\textbf{完成初版基于大模型的Live2D数字人驱动方案验证}，实现离线可用素材库生成并保障Demo效果；
\item 多语言大模型构建和落地 (2024.3-2024.8)：(1) 负责小艺多语言大模型全流程构建（高效多语言词表扩充、多语言预训练数据清洗、多语言训练配比，多语言后训练等），牵引诺亚实现盘古多语言模型能力显著提升；(2) 基于mLLM显著提升小艺翻译体验，探索1.5B-LLM端侧大模型结合语音噪音数据增强策略，实现通话翻译场景中英双向性能平均提升7个COMET值，\textbf{体感测试显著优于三星手机翻译30\%+}；7B-LLM中英互译能力在开源翻译测试集上持平GPT-4性能，显著优于线上最优模型；(3) 协助整个小艺翻译团队梳理整个技术方案和业务分工；
\end{itemize}

\datedsubsection{\textbf{阶跃星辰}，北京}{2023.11 -- 2024.3}
\text{算法专家，主要负责语言大模型预训练工作以及前沿技术探索}
\begin{itemize}
\item Step2近万亿大模型预训练项目技术骨干，深度参与MoE结构和训练算法策略验证；主导FP8-LM训练技术探索和落地（包括预训练和后训练阶段），推动FP8相关代码库完善；
\item 负责大模型长序列策略探索，验证各类主流LongContext策略，支撑128K上下文落地；
\end{itemize}

\datedsubsection{\textbf{腾讯-AI Lab}，深圳}{2021.9 -- 2023.11}
\role{技术级别: T11, 自然语言算法平台组，1次5星, 2次SEVP卓越个人}
\text{主要负责腾讯翻译训练平台搭建、交互翻译模型建设、个性化机器翻译以及先进技术研究}
\begin{itemize}
  \item 负责腾讯翻译分布式训练平台搭建和优化，自动化训练流程搭建（包括高质量翻译数据过滤和领域识别），交互翻译多合一模型(具备自动翻译、约束解码和翻译记忆融合功能)构建与优化，个性化机器翻译策略研发与落地，参与翻译商业化项目；协助管理自然语言算法平台组翻译方向人员；
  \item 开展自然语言处理相关研究，提出联邦机器翻译,交互机器翻译开源平台, Simple and Scalable Nearest Neighbor Machine Translation (SK-MT)等新方法，其中\textbf{SK-MT成为腾讯翻译的个性化机器翻译核心策略并完成上线}；另外探索模型鲁棒性相关方法，包括理解词级别文本攻击和通过Nearest Neighbor Calibration提升LLM的In-Context Learning能力; 
  \item 开展多模态相关机器翻译探索，包括Speech-to-Speech Translation和Automatic Song Translation等方向，取得IWSLT2023比赛中\textbf{端到端语音到语音翻译子任务第一名}；
  \item 参与混元大模型项目，负责多语言预训练数据清洗，LLM的翻译能力和RLHF相关调优；
\end{itemize}

\datedsubsection{\textbf{阿里巴巴-达摩院}，杭州}{2019.7 -- 2021.8}
\role{算法专家，语言技术实验室，阿里星，1次375绩效}
% {主管：陈博兴、谢军}
\text{主要负责阿里翻译基础通用模型建设、语音翻译、对外商业化以及先进技术研究}
\begin{itemize}
  \item 负责阿里云通用翻译服务以及商业化项目，技术上完成阿里通用翻译自动化链路构建、通用模型质量优化（中英互译内部评测达到Top1）、通用语向扩展（15->214）、线上模型解码加速（2-3x)，以及商业化领域模型定制和交付;
  \item 开展自然语言处理相关研究，提出高效的多语言机器翻译训练策略，融合预训练语言模型的机器翻译，Iterative Domain-Repaired Back-Translation，端到端语音翻译模型，以及初步探索个性化机器翻译策略。这些方法已用于大幅度提升线上多语言翻译模型性能和领域模型定制效果。另外也探索可变长的文本对抗攻击和任务型对话系统作为自然语言生成等新方法;
  \item 受邀在DataFunTalk 2020年终大会上分享“\textbf{阿里多语言翻译模型的前沿探索及技术实践}”的主题报告;
\end{itemize}

\datedsubsection{\textbf{微软亚洲研究院}，北京}{2019年2月 -- 2019年6月}
\role{研究实习生，自然语言计算组，(联合培养博士生项目)}{指导老师: 刘树杰}
\textit{自然语言处理相关技术研究}

\datedsubsection{\textbf{微软雷德蒙德研究院}，西雅图，美国}{2018年7月 -- 2019年1月}
\role{Research Intern，Deep Learning Group,}{指导老师: Xiujun Li, Jianfeng Gao}
\textit{对话系统相关技术研究}
\begin{itemize}
  \item 面向资源受限的任务型对话场景，提出Budgeted Policy Learning来显著提升对话系统的成功率；
\end{itemize}

\datedsubsection{\textbf{微软亚洲研究院}，北京}{2015年7月 -- 2018年7月}
\role{研究实习生，自然语言计算组，(联合培养博士生项目)}{指导老师: 李沐}
\textit{机器翻译、深度学习、自然语言生成等相关技术研究}
\begin{itemize}
  \item 开展神经机器翻译相关研究，提出Iterative Back-Translation, Target-bidirectional Agreement, Coarse-to-Fine Learning, Bidirectional GANs等训练算法，针对无监督机器翻译设计SMT as Posterior Regularization方法，另外将NMT相关技术应用于Dependency Parsing和Text Style Transfer任务中;
  \item 支持微软内部项目开发(SmartFlow Toolkit、Writing Intelligence和Babel Project):
  \item[-] \emph{SmartFlow Toolkit (2016)}:  开发C\#版本的深度学习工具，实现OP计算、计算图调度，内部管理等功能，并服务于组内各类模型开发（包括R-NET模型、对话/QA相关的NLG模型等），\textbf{部分模型已应用于微软小冰对话生成和必应产品中};
  \item[-] \emph{Writing Intelligence Project (2017)}: 该项目想借助深度学习技术去使写作更便捷，其中定义了句子补全、基于关键词的句子生成、语法检查、下一句话推荐这四个功能。主要负责句子补全功能的模型开发以及实现;
  \item[-] \emph{Babel Project (2017-2018)}: 微软内部多个小组一起合作，探索在大规模语料下当前神经机器翻译的极限。负责将Iterative Back-Translation和Target-bidirectional Agreement两个技术应用于该项目，显著提升模型的性能，\textbf{并达到人类可比水平};
\end{itemize}

\datedsubsection{\textbf{微软亚洲研究院}，北京}{2013年7月 -- 2014年6月}
\role{研究实习生，自然语言计算组}{指导老师: 李沐}
\textit{自然语言处理相关技术研究}
\begin{itemize}
  \item 熟悉和复现各类基础NLP模型，将word2vec技术用于篇章关系分析中；
\end{itemize}

\section{\faFilesO\ 主要论文和技术报告}

\begin{enumerate}

\item \normalsize \textbf{Safe Delta: Consistently Preserving Safety when Fine-Tuning LLMs on Diverse Datasets}\\
\small - Ning Lu1, Shengcai Liu, Jiahao Wu, Weiyu Chen, \textbf{Zhirui Zhang}, Yew-Soon Ong, Qi Wang, Ke Tang. In \textit{ICML}'2025

\item \normalsize \textbf{Lost in the Source Language: How Large Language Models Evaluate the Quality of Machine Translation}\\
\small - Xu Huang, \textbf{Zhirui Zhang}, Xiang Geng, Yichao Du, Jiajun Chen, Shujian Huang. In \textit{ACL}'2024

\item \normalsize \textbf{Document-Level Machine Translation with Large Language Models}\\ 
\small - Longyue Wang*, Chenyang Lyu*, Tianbo Ji*, \textbf{Zhirui Zhang*}, Dian Yu, Shuming Shi, Zhaopeng Tu. In \textit{EMNLP}'2023.

\item \normalsize \textbf{Nearest Neighbor Machine Translation is Meta-Optimizer on Output Projection Layer}\\ 
\small - Ruize Gao,  \textbf{Zhirui Zhang}, Yichao Du, Lemao Liu, Rui Wang. In \textit{EMNLP}'2023.

\item \normalsize \textbf{Fairness-guided Few-shot Prompting for Large Language Models}\\ 
\small - Huan Ma, Changqing Zhang, Yatao Bian, Lemao Liu, \textbf{Zhirui Zhang}, Peilin Zhao, Shu Zhang, Huazhu Fu, Qinghua Hu, Bingzhe Wu. In \textit{NeurIPS}'2023.

\item \normalsize \textbf{The \textsc{MineTrans} Systems for IWSLT 2023 Offline Speech Translation and Speech-to-Speech Translation Tasks}\\ 
\small - Yichao Du, Zhengsheng Guo, Jinchuan Tian, \textbf{Zhirui Zhang}, Xing Wang, Jianwei Yu, Zhaopeng Tu, Tong Xu and Enhong Chen. In \textit{IWSLT}'2023.

\item \normalsize \textbf{E-NER: Evidential Deep Learning for Trustworthy Named Entity Recognition}\\
\small - Zhen Zhang, Mengting Hu, Shiwan Zhao, Minlie Huang, Haotian Wang, Lemao Liu, \textbf{Zhirui Zhang}, Zhe Liu and Bingzhe Wu.  In \textit{ACL}'2023.

\item \normalsize \textbf{Federated Nearest Neighbor Machine Translation}\\
\small - Yichao Du, \textbf{Zhirui Zhang}, Bingzhe Wu, Lemao Liu, Tong Xu and Enhong Chen. In \textit{ICLR}'2023.

\item \normalsize \textbf{Simple and Scalable Nearest Neighbor Machine Translation}\\
\small - Yuhan Dai*, \textbf{Zhirui Zhang*}, Qiuzhi Liu, Qu Cui, Weihua Li, Yichao Du and Tong Xu. In \textit{ICLR}'2023.

\item \normalsize \textbf{Task-Oriented Dialogue System as Natural Language Generation}\\
\small - Weizhi Wang, \textbf{Zhirui Zhang}, Junliang Guo, Yinpei Dai, Boxing Chen and Weihua Luo. In \textit{SIGIR}'2022.

\item \normalsize \textbf{Automatic Song Translation for Tonal Languages}\\
\small - Fenfei Guo, Chen Zhang, \textbf{Zhirui Zhang}, Qixin He, Kejun Zhang, Jun Xie and Jordan Lee Boyd-Graber. In \textit{ACL}'2022.

\item \normalsize \textbf{Adaptive Nearest Neighbor Machine Translation}\\
\small - Xin Zheng, \textbf{Zhirui Zhang}, Junliang Guo, Shujian Huang, Boxing Chen, Weihua Luo and Jiajun Chen. In \textit{ACL-IJCNLP}'2021.

\item \normalsize \textbf{Incorporating BERT into Parallel Sequence Decoding with Adapters}\\
\small - Junliang Guo, \textbf{Zhirui Zhang}, Linli Xu, Hao-Ran Wei, Boxing Chen and Enhong Chen. In \textit{NeurIPS}'2020

\item \normalsize \textbf{Cross-lingual Pre-training Based Transfer for Zero-shot Neural Machine Translation} \\
\small - Baijun Ji, \textbf{Zhirui Zhang}, Xiangyu Duan, Min Zhang, Boxing Chen and Weihua Luo. In \textit{AAAI}'2020.

\item \normalsize \textbf{Budgeted Policy Learning for Task-Oriented Dialogue Systems} \\
\small - \textbf{Zhirui Zhang}, Xiujun Li, Jianfeng Gao and Enhong Chen. In \textit{ACL}'2019.

\item \normalsize \textbf{Regularizing Neural Machine Translation by Target-bidirectional Agreement} \\
\small - \textbf{Zhirui Zhang}, Shuangzhi Wu, Shujie Liu, Mu Li, Ming Zhou and Tong Xu. In \textit{AAAI}'2019.

\item \normalsize \textbf{Bidirectional Generative Adversarial Networks for Neural Machine Translation} \\
\small - \textbf{Zhirui Zhang}, Shujie Liu, Mu Li, Ming Zhou and Enhong Chen. In \textit{CoNLL}'2018.

\item \normalsize \textbf{Joint Training for Neural Machine Translation Models with Monolingual Data} \\
\small - \textbf{Zhirui Zhang}, Shujie Liu, Mu Li, Ming Zhou and Enhong Chen. In \textit{AAAI}'2018.

\item \normalsize \textcolor{red}{\textbf{Achieving Human Parity on Automatic Chinese to English News Translation}} \\
\small - Hany Hassan, Anthony Aue, Chang Chen, Vishal Chowdhary, Jonathan Clark, Christian Federmann, Xuedong Huang, Marcin Junczys-Dowmunt, William Lewis, Mu Li, Shujie Liu, Tie-Yan Liu, Renqian Luo, Arul Menezes, Tao Qin, Frank Seide, Xu Tan, Fei Tian, Lijun Wu, Shuangzhi Wu, Yingce Xia, Dongdong Zhang, \textbf{Zhirui Zhang} and Ming Zhou. \textit{Pre-print}, https://arxiv.org/abs/1803.05567 

\item \normalsize \textbf{Style Transfer as Unsupervised Machine Translation} \\
\small - \textbf{Zhirui Zhang*}, Shuo Ren*, Shujie Liu, Jianyong Wang, Peng Chen, Mu Li, Ming Zhou and Enhong Chen. \textit{Pre-print}, https://arxiv.org/abs/1808.07894

\end{enumerate}

\section{\faCogs\ 主要技能及学术服务}
% increase linespacing [parsep=0.5ex]
\begin{itemize}
  \item 编程技能: Python/C++/C\#/Java/CUDA/Cudnn/Tensorflow/Pytorch
  \item 学术服务: 担任多个NLP/ML会议与期刊长期审稿人，包括ACL、EMNLP、NAACL、AAAI、IJCAI、NeurIPS、ICML、ICLR，COLR等，担任过ACL领域主席;
\end{itemize}

\section{\faHeartO\ 主要获奖情况}
\datedline{\textit{国家奖学金}}{2013}
\datedline{\textit{谷歌奖学金}}{2013}
\datedline{\textit{微软亚洲研究院实习生“明日之星”}}{2019}

%% Reference
%\newpage
%\bibliographystyle{IEEETran}
%\bibliography{mycite}
\end{document}
